\sectionkicker{Source-backed technical notes}
\subsection*{もう一つの生成系譜――DDPMからscore-SDEへ: Technical Notes}
本欄は記事本文で圧縮した一次資料上の情報を、比較・再検証しやすい形で整理したものである。時系列、確認済みの主張、留意点、一次資料URLを掲載する。
\smallskip
\noindent{\small\textit{Technical Notesでは、一次資料から確認した公開・提供・研究上の事実を記載する。数値・能力評価や提供元・著者による比較表現は、独立再現済みの結果ではなく帰属claimとして扱う。}}\par
\medskip
\subsection*{Theme at a glance}
\addcontentsline{toc}{subsection}{Theme at a glance}
\begin{center}
\footnotesize
\begin{tabularx}{\linewidth}{@{}p{0.20\linewidth}p{0.10\linewidth}p{0.14\linewidth}X@{}}
\toprule
資料 & 位置づけ & 種別 & 時系列 \\
\midrule
Denoising Diffusion Probabilistic Models & 主要資料 & 論文 & 2020-06-19 \\
Score-Based Generative Modeling through Stochastic Differential Equations & 主要資料 & 論文 & 2020-11-26 \\
\bottomrule
\end{tabularx}
\end{center}
\normalsize

\begin{technicalnote}{Denoising Diffusion Probabilistic Models}{主要資料}
\begingroup
% reader-facing Technical Notes break policy
\widowpenalty=10000
\clubpenalty=10000
\displaywidowpenalty=10000
\begin{tabularx}{\linewidth}{@{}>{\bfseries}p{0.22\linewidth}X@{}}
組織 & UC Berkeley / authors \\
種別 & 論文 \\
時系列 & 2020-06-19 \\
\end{tabularx}
\smallskip
\begin{minipage}{\linewidth}
% reader-facing Technical Notes event-only fact/source tail group
{\bfseries 一次資料から整理したtechnical points}
\begin{itemize}[leftmargin=1.5em,itemsep=0.35em]
\item \textbf{一次情報で確認できる事実}: UC Berkeley / authorsの選定済み一次資料では、DDPMはdiffusion probabilistic latent-variable modelをweighted variational boundでtrainingし、denoising score matching with Langevin dynamicsとの接続を用いる。 生成過程はprogressive lossy decompressionとして解釈でき、autoregressive decodingの一般化という見方を提示する。 CIFAR-10やLSUNでのsample-quality指標は著者実験の結果として扱う。
\end{itemize}
\begin{samepage}
{\bfseries 一次資料}
\begingroup\sloppy
\Urlmuskip=0mu plus 2mu\relax
\begin{itemize}[leftmargin=1.5em,itemsep=0.25em]
\item {\scriptsize\url{http://arxiv.org/abs/2006.11239v2}}
\end{itemize}
\endgroup
\end{samepage}
\end{minipage}
\endgroup
\end{technicalnote}
\medskip

\begin{technicalnote}{Score-Based Generative Modeling through Stochastic Differential Equations}{主要資料}
\begingroup
% reader-facing Technical Notes break policy
\widowpenalty=10000
\clubpenalty=10000
\displaywidowpenalty=10000
\begin{tabularx}{\linewidth}{@{}>{\bfseries}p{0.22\linewidth}X@{}}
組織 & Stanford / authors \\
種別 & 論文 \\
時系列 & 2020-11-26 \\
\end{tabularx}
\smallskip
\begin{minipage}{\linewidth}
% reader-facing Technical Notes event-only fact/source tail group
{\bfseries 一次資料から整理したtechnical points}
\begin{minipage}{\linewidth}
% reader-facing Technical Notes coherent tail group
\begin{itemize}[leftmargin=1.5em,itemsep=0.35em]
\item \textbf{一次情報で確認できる事実}: Stanford / authorsの選定済み一次資料では、Score-SDEはforward SDEでdata distributionへ徐々にnoiseを加えてknown priorへ移し、time-dependent scoreに依存するreverse-time SDEでpriorからdataへ戻す。 score networkとnumerical SDE solverによるsamplingに加え、predictor-corrector frameworkとequivalent neural ODEを提示し、score-based generative modelingとdiffusion probabilistic modelingを同一frameworkに収める。 exact-likelihood計算やhigh-resolution generation等の結果は論文条件下の著者実験として扱う。
\end{itemize}
\begin{samepage}
{\bfseries 一次資料}
\begingroup\sloppy
\Urlmuskip=0mu plus 2mu\relax
\begin{itemize}[leftmargin=1.5em,itemsep=0.25em]
\item {\scriptsize\url{http://arxiv.org/abs/2011.13456v2}}
\end{itemize}
\endgroup
\end{samepage}
\end{minipage}
\end{minipage}
\endgroup
\end{technicalnote}
\medskip
